\scanpage{23}
\fbox{Dans le cas d'un \uline{motif} $E$,}

On va exprimer $\delta(E)$ et $\chi(E)$ en termes des \deleted{\illegible{1}} composantes virtuelles des $\mathbb R f_!(E)$ (utilisant conj. de Weil), on trouve
\[
\left\{\begin{aligned}
\chi(E)&=\sum(-1)^i\beta_i(E)\\
\delta(E)&=\prod_i\delta\bigl(h^i_!(E)\bigr)^{(-1)^i}
\end{aligned}\right.
\]
\margin{$\beta_i(E)$ les \unsure{nbs|nb\textsuperscript{s}} de Betti virtuels à coeff. dans $E$}

Or $h^i_!(E)$ donne $\beta_i$ valeurs propres de val. absolue $p^{i/2}$, donc
\[
\deleted{\unsure{\delta(\beta_i)=|\delta(\beta)=}}\delta\bigl(h^i_!(E)\bigr)=\varepsilon_i\,p^{\frac{i\beta_i}{2}}
\]
\margin{(NB. $i\beta_i/2$ tjrs entier)}

où $\deleted{\unsure{\delta|\varepsilon}}\varepsilon_i$ est défini comme un signe de déterminant virtuel, et $\varepsilon_i$ est $\deleted{\unsure{=|\varepsilon}}+1$ si $i$ est impair. On trouve alors
\[
\left\{\begin{aligned}
\delta(E)&=\deleted{\unsure{\varepsilon(E)p^{\frac12\sum i\beta_i}|\varepsilon(E)p^{\sum i\beta_i}}}\varepsilon(E)\,p^{\frac12\sum_i(-1)^i\,i\,\beta_i(E)}\\
\varepsilon(E)&=\prod_i\varepsilon_i(E)
\end{aligned}\right.
\]
\margin{$\varepsilon(E)$ et les $\varepsilon_i(E)$ sont $\pm1$, $\varepsilon_i(E)=+1$ si $i$ impair.}
% archive sheet number: 22
